\section{Paths and formula depth}\label{sec:formula}

The state sum on a path is iterated matrix multiplication. With $d\ge2$
vertices and $q$ states per edge, it is the $(1,1)$ entry of a product
of $d$ generic $q\times q$ matrices, using only the first row and last
column at the endpoints. An ABP realizes this product with one source,
$d-1$ layers of $q$ states, and one sink. In particular, taking $d=q=n$
gives a family in $\VBP$. To obtain a formula, bisect the path and sum
over the $q$ possible labels on the middle edge. Each summand computes
two shorter entries, so recursive bisection gives size
$q^{O(\log d)}$ for $q\ge2$. On short paths this exponent is optimal,
even for degree-component border formulas. For the monotone and
set-multilinear versions of iterated matrix multiplication, see
Snir~\cite{Snir} and Tavenas, Limaye and Srinivasan~\cite{LST}.

\begin{theorem}[Short trees]\label{thm:full-flat-graph}
Let $G$ be any tree on $d$ vertices and let $w=\log q$.
For $2\le d\le w^{1/3}$ and all sufficiently large $q$, the
degree-component border formula size of $Z_{G,q}$ is at least
$q^{c\log d}$ over every field. Moreover, if $Z_{G,q}$ lies in the
Zariski closure of the degree-$d$ components of outputs of depth-$h$
binary arithmetic circuits over the algebraic closure, then
$h\ge c' w\log d$. The constants $c,c'>0$ are absolute.
\end{theorem}

A cut between two consecutive path vertices crosses just one edge.
To recover a factor of $\log d$ in the lower bound we use cuts at
several scales in a single placement. The cuts depend on how a
formula splits the vertex groups among its products. That splitting
defines a binary tree on the groups, whatever the shape of the original
graph $G$. We first describe the information shared among descendants
of a split, then choose sites where the price of this information is
controlled by the formula depth.

\subsection{The coefficient tensor of a formula}

Work with the $d$ slots of Section~\ref{sec:cuts}, keeping each
vertex table at a single site for now. For a set $B$ of slots, write
$\Delta_B(2^\ell)$ for the diagonal tensor supplying one $\ell$-bit
label to every slot in $B$. When $\ell=0$ this is the tensor of all
local units. In a rooted tree whose leaves are the slots, let $B_s$
be the set of leaves below a node $s$.

\begin{lemma}[Coefficient extraction at the product gates]\label{lem:formula-splits}
For a binary formula of depth $h$ with output $f$, the tensor $A(f)$
of coefficients using one variable from each slot is an ordinary sum
of at most $4^d d!(h+1)^{d-1}$ tensors. Each summand is a local linear
image, with one map at each slot, of
\[
                \bigotimes_{s\ \mathrm{internal}}\Delta_{B_s}(2^{\ell_s})
\]
for some ordered full binary tree on the slots and nonnegative integers
$\ell_s$ satisfying
\[
                \sum_{s:\,v\in B_s}\ell_s\le h
                \qquad\text{at every leaf }v .
\]
\end{lemma}

\begin{proof}
For a subformula rooted at an occurrence $u$, extract the coefficient
table on a nonempty set $B$ of slots. Its entries are the coefficients
of monomials containing one chosen variable from each slot of $B$
and no other variables. For a singleton $B$ this is the vector of
linear coefficients of that slot. For $|B|\ge2$, expand down the
formula to the first product occurrence $p$ at which both children
supply at least one variable. Such a $p$ can be $u$ itself.
At each intervening product all variables come from just one child,
so the unused child contributes its constant term. At a subtraction
the right child contributes a minus sign. Denote the product of these
path scalars by $c(u,p)$.

Consequently the table on $B$ is the sum over $p$ and over ordered
partitions $B=B_1\sqcup B_2$ with $B_1,B_2$ nonempty. Each term is
$c(u,p)$ times the tensor product of the extracted tables on $B_1$
and $B_2$ at the respective children of $p$, with slots returned to
their original order. This is the coefficient rule at the gates.
An affine leaf contributes zero to a table with $|B|\ge2$.

Apply the rule recursively from the output. Group the resulting terms
by the ordered binary tree $\tau$ of nontrivial splits, the bijection
from its leaves to the $d$ slots, and, at each split $s$, the distance
$\ell_s$ from the root of its subformula to the chosen product occurrence.
There are fewer than $4^d$ possible shapes, $d!$ labelings and at most
$(h+1)^{d-1}$ lists of distances. Along any one leaf's chain of
calls, the traversed paths occupy disjoint portions of a path in the
formula. Their lengths therefore sum to at most $h$.

Fix such a group. At split $s$, take an independent $\ell_s$-bit
address, one bit for each downward choice of child on the way to $p$.
Supply this address to all leaves in $B_s$. A leaf decodes its
ancestor addresses in order. At each decoded product it knows from
$\tau$ which child starts the next call. It can thus return the
linear coefficient vector at its final call, or zero if an address
fails to reach a product at a prescribed split.
For each internal node designate one descendant leaf to multiply its
vector by the scalar $c(u,p)$ for that node. Each leaf's rule on
basis labels defines a linear map. Summing over the independent
addresses in the product of diagonals now sums over exactly the
consistent choices of products, each with its scalar appearing once.
The resulting local image is the required group of terms.
\end{proof}

The price of a shared address depends on where its descendant leaves
are placed. The next estimate chooses those sites separately for each
tree $\tau$. The terminal positions of Section~\ref{sec:cuts} stay
fixed, so bounds obtained for different $\tau$ can eventually be added
using Equation~\eqref{eq:ordinary}.

\subsection{Placing the slots at several scales}

Recall that $a_v$ is the sum of the unit terminal vectors at vertex
$v$, and $a_C=\sum_{v\in C}a_v$. Connectivity of $G$ gives
$\|a_C\|\ge1$ for every nonempty proper vertex set $C$, by
Equation~\eqref{eq:cut}. A displacement parallel to $a_C$ at the
vertices of $C$ shortens the total distance to their terminals to
first order. We choose one displacement per scale so that the
shortenings add while each descendant set remains narrow.

\begin{lemma}[Placement of the leaves]\label{lem:formula-placement}
For any ordered full binary tree $\tau$ on the $d$ slots of a connected
graph $G$ with $d\ge2$, there are vectors $z_v$ in the terminal space
with $\|z_v\|\le1$ such that
\[
       \sum_v\langle a_v,z_v\rangle\ge c_0\log d,
       \qquad
       \delta(\{z_v:v\in B_s\})\le C_0\,\frac{|B_s|}{d}
       \quad\text{for every node }s .
\]
Here $\delta$ is the normalized Gaussian mean width of
Theorem~\ref{thm:width}, and $c_0,C_0>0$ are absolute constants.
\end{lemma}

\begin{proof}
Read the leaves in the order of $\tau$. For each power of two $k\le d/2$,
the maximal subtrees having at most $k$ leaves partition that order into
consecutive blocks. Collect blocks from left to right, closing a cluster
as soon as it has at least $k$ leaves. Every cluster has fewer than
$2k$ leaves, and only the last can have fewer than $k$. Thus each is a
proper subset of the $d$ slots and there are at least $d/(2k)$ clusters.
Figure~\ref{fig:cluster} shows how the same split is seen at two scales.

\begin{figure}[!hbp]
\centering
\begin{tikzpicture}[x=11mm,y=9mm,font=\small,
  branch/.style={draw=ink!80,line width=.65pt},
  marked/.style={draw=sea,line width=1pt},
  cluster/.style={draw=ink!55,rounded corners=2pt,fill=white},
  meets/.style={draw=sea,rounded corners=2pt,fill=sea!12},
  block/.style={|-|,draw=ink,line width=.7pt}]
\foreach \i in {1,...,8}
  \coordinate (leaf\i) at (\i,0);
\coordinate (root) at (4.1,3.2);
\coordinate (left) at (2,2.05);
\coordinate (right) at (5.7,2.3);
\coordinate (leftpair) at (2.5,.85);
\coordinate (s) at (6.5,1.65);
\coordinate (pairfive) at (5.5,.75);
\coordinate (pairseven) at (7.5,.75);
\draw[branch] (left)--(root)--(right)
              (leaf1)--(left)--(leftpair)
              (leaf2)--(leftpair)--(leaf3)
              (leaf4)--(right)--(s);
\draw[marked] (pairfive)--(s)--(pairseven)
              (leaf5)--(pairfive)--(leaf6)
              (leaf7)--(pairseven)--(leaf8);
\foreach \n in {root,left,right,leftpair} \fill[ink] (\n) circle (1.4pt);
\foreach \n in {s,pairfive,pairseven} \fill[sea] (\n) circle (1.6pt);
\node[right=3pt,text=sea] at (s) {$s$};
\foreach \i in {1,...,8} {
  \fill[ink] (leaf\i) circle (1.4pt);
  \node[below=1pt] at (leaf\i) {\(\i\)};
  \draw[densely dotted,ink!45] (\i,-.52)--(\i,-1.92);
}
\foreach \k/\y in {2/-.9,4/-1.72}
  \node[anchor=east] at (.25,\y) {\(k=\k\)};
\draw[cluster] (.55,-.65) rectangle (3.45,-1.15);
\foreach \l/\r in {4/6,7/8}
  \draw[meets] (\l-.45,-.65) rectangle (\r+.45,-1.15);
\draw[cluster] (.55,-1.47) rectangle (4.45,-1.97);
\draw[meets] (4.55,-1.47) rectangle (8.45,-1.97);
\foreach \l/\r in {1/1,2/3,4/4,5/6,7/8}
  \draw[block] (\l-.30,-.9)--(\r+.30,-.9);
\foreach \l/\r in {1/3,4/4,5/8}
  \draw[block] (\l-.30,-1.72)--(\r+.30,-1.72);
\end{tikzpicture}
\caption{An ordered tree of formula splits, with slots numbered in leaf
order. The bars show the maximal subtrees of at most \(k\) leaves, and
the boxes group consecutive bars into clusters. The highlighted subtree
has \(B_s=\{5,6,7,8\}\). It meets two clusters when \(k=2\) and lies in
one when \(k=4\), as the shaded boxes show. All four leaves therefore
receive the same displacement at the larger scale.}
\label{fig:cluster}
\end{figure}


Write $C_k(v)$ for the cluster of $v$. Displace every member of a
cluster in the same unit direction $a_C/\|a_C\|$ with weight $k/d$.
The combined displacement is
\[
            z_v=\sum_k\frac{k}{d}\frac{a_{C_k(v)}}{\|a_{C_k(v)}\|}.
\]
The weights sum to at most one. At a given scale, the inner products
within a cluster sum to $(k/d)\|a_C\|$. Thus
\[
       \sum_v\langle a_v,z_v\rangle
         =\sum_k\frac{k}{d}\sum_{C\text{ at scale }k}\|a_C\|
         \ge \tfrac12\lfloor\log d\rfloor .
\]

For the width bound, fix a node $s$ with $b=|B_s|$. At each scale
$k\ge b$ its entire subtree belongs to a single maximal block. Its
displacement at that scale is a common translation, which leaves its
width unchanged. At a scale $k<b$, the interval $B_s$ meets at most
$2+b/k$ clusters. All except possibly the two end clusters lie wholly
in the interval and have at least $k$ leaves. There are therefore at
most $2+b/k$ unit displacement directions on $B_s$ at this scale.

The set of combined displacements on $B_s$ is contained, up to
translation, in the Minkowski sum of these sets of directions scaled
by $k/d$. In the Gaussian formula, taking a maximum of a linear
functional gives monotonicity under inclusion and additivity on
Minkowski sums. Applying Equation~\eqref{eq:width-ball} to each set
of unit directions gives
\[
       \delta(\{z_v:v\in B_s\})
          \le C\sum_{k<b}\frac{k}{d}\sqrt{\log(2+b/k)}
          \le C_0\,\frac{b}{d}.
\]
At the $j$th dyadic scale below $b$, the summand is at most a
constant times $(b/d)2^{-j}\sqrt{j+2}$, whose sum converges.
\end{proof}

\subsection{From depth to size}

\begin{proof}[Proof of Theorem~\ref{thm:full-flat-graph}]
Let $e=d-1$ be the number of edges of $G$. Its split coefficient
tensor has terminal price $2ew$. Take a formula of depth $h$, and one
summand in Lemma~\ref{lem:formula-splits}. For this summand choose
the vectors $z_v$ of Lemma~\ref{lem:formula-placement} and perform the
map at slot $v$ at the site $\lambda z_v$, with $\lambda\ge0$.
Share each address among its descendants at those sites, then send the
$\deg(v)$ output indices of each slot to their terminals by $q$-state
wires. The resulting upper bound on the terminal price is
\begin{equation}\label{eq:full-flat-profile}
 \begin{aligned}
   2ew
     &-w\lambda\sum_v\langle a_v,z_v\rangle
      +\tfrac12 w\lambda^2\sum_v\deg(v)\|z_v\|^2\\
     &+\lambda\sum_s\ell_s\,\delta(\{z_v:v\in B_s\}).
 \end{aligned}
\end{equation}
The first line uses Equation~\eqref{eq:move} for the output wires.
The second is the exact price of the independent address labels.
Its coefficient of $\lambda$ is at most $C_0h$, since
\[
        \sum_s\ell_s\frac{|B_s|}{d}
            =\frac1d\sum_v\sum_{s:\,v\in B_s}\ell_s\le h .
\]
The first-order shortening in the first line is at least
$c_0w\lambda\log d$, and the quadratic term is at most $ew\lambda^2$.
If $h\le c_1w\log d$ for sufficiently small $c_1$, choosing
$\lambda=c_2(\log d)/d$ with $c_2>0$ sufficiently small bounds
the price of each summand by
\[
                  2ew-c_3\,\frac{w(\log d)^2}{d}
\]
for an absolute $c_3>0$.

All summands now have the same $2e$ terminal sites. By
Equations~\eqref{eq:ordinary} and~\eqref{eq:width-ball}, adding the
summands of Lemma~\ref{lem:formula-splits} costs a premium of at most
\[
                O\bigl(d\sqrt{\log(2d)}\,\log(d(h+2))\bigr).
\]
The ratio of this premium to the saving is $O(w^{-1/3}\log w)$,
uniformly for $2\le d\le w^{1/3}$ and $h\le c_1w\log d$. Thus the
coefficient tensors of all such formulas lie in a fixed price
sublevel strictly below $2ew$ for all sufficiently large $w$.
The sublevel is Zariski closed, and extraction of one variable per
slot is a linear map on degree-$d$ components. The graph polynomial
is therefore excluded even from the degree-component border of the
depth-$h$ formulas. Unfolding a binary circuit into a formula
preserves its depth and its output, giving the claimed circuit depth
bound as well.

Finally, a formula of size $m$ can be balanced to depth
$O(\log(m+2))$ over every field~\cite{Brent}. For the division-free
construction, push subtraction signs to the affine leaves and select
a subformula $g$ containing between a third and two thirds of the
leaves. The output equals $ag+b$, where $a$ is the product of the
siblings at the multiplication gates on the path from $g$ to the
root, and $b$ is the original expression evaluated with $g=0$.
Each of $a,b,g$ uses at most two thirds of the original leaves plus
one. Recursion down to constant size gives the desired depth.
Balancing preserves the polynomial before closure, so the depth
bound yields the degree-component border formula size in the theorem.
\end{proof}

For a path, Theorem~\ref{thm:full-flat-graph} and the bisection formula
give matching exponents. The estimate extends to
$2\le d\le(\log q)^B$ for any fixed $B>1/3$, with lower bound
$q^{\Omega_B(\log d)}$ on every tree in this range. When
$d>(\log q)^{1/3}$, project to a subtree on
$\lfloor(\log q)^{1/3}\rfloor$ vertices. The logarithm of this
subtree size is $\Omega_B(\log d)$, and the quadratic cost of the
border formula projection in Lemma~\ref{lem:projection} only changes
the implied constant in the exponent.

Taking $q=n$ and projecting the $n$-vertex path to its first
$\lfloor(\log n)^{1/3}\rfloor$ vertices gives degree-component
border formula size $n^{\Omega(\log\log n)}$. Its layered ABP has
polynomial size, so this proves the path separation in
Theorem~\ref{thm:main}.

The same lower bound transfers to the order-$n$ determinant. Take
$q=\lfloor\sqrt n\rfloor$ and $d=\lfloor(\log q)^{1/3}\rfloor$.
Let $A$ be the weighted adjacency matrix of the layered ABP in
topological order. It is strictly upper triangular, so
$\det(I-A)=1$. The source--sink entry of
$(I-A)^{-1}=\sum_{j\ge0}A^j$ is the path polynomial. By the
cofactor formula it is also a signed minor of $I-A$ of order
$(d-1)q+1\le n$~\cite{Valiant,MalodPortier}. Absorb the sign into
one row and pad with an identity block to obtain an affine projection
of the order-$n$ determinant. Lemma~\ref{lem:projection} transfers
the short-path lower bound and completes the proof of
Theorem~\ref{thm:main}.
